\usepackage[utf8]{inputenc}
% Para cocina
\usepackage{xcookybooky}
% Idioma
\usepackage[english, spanish]{babel}
% Indice, o urls
\usepackage[bookmarks = true, colorlinks=true, linkcolor = black, citecolor = black, menucolor = black, urlcolor = blue]{hyperref}

% Título
\newcommand{\titulo}{Recopilación De Recetas}
\newcommand{\yo}{Ex Xavi Salvaje}

% *** COLOR ***
\usepackage[table]{xcolor}
\definecolor{azul}{RGB}{0,160,255}
\definecolor{azul2}{RGB}{7,74,163}
\definecolor{azul3}{RGB}{0,89,140}
\definecolor{alizarin}{rgb}{0.82, 0.1, 0.26}
\definecolor{gray}{RGB}{242,242,242}
\definecolor{grayblack}{RGB}{50,50,50}


\usepackage{hyperref}
\hypersetup{
	pdftitle={\titulo},
	pdfsubject={\yo},
	pdfauthor={\yo},
	colorlinks=true, 
	linkcolor=azul, 
	citecolor=cyan, 
	urlcolor=azul, 
	linktocpage,
	pdfkeywords={recetas, cocina}
}
% Paquetes de imagenes
\usepackage{graphicx}
\usepackage{wrapfig}
% Ruta de imágenes
\graphicspath{ {pix/} }
% Lenguaje
\selectlanguage{spanish}
% Interlineado
\usepackage{setspace}
% Para paginaciones
\usepackage{fancyhdr}
% Limpia cabecera y pie de página
\fancyhead{}
\fancyfoot{}
% Setea número hacia el lado derecho del pie de página
\fancyfoot[R]{\thepage}
% Configura xcookybooky
\setHeadlines{
	inghead = Ingredientes,
	prephead = Preparación,
	hinthead = Tips
}

\title{\titulo}
\author{\yo}